% GENERATED FILE. Edit canonical lessons, then run npm run generate:books.
% canonical-source-hash: fnv1a64:c5cd075bf1b65a00
% canonical-lessons: ZH-C44-che, ZH-W44-che, ZH-C44-fandian, ZH-C44-yiyuan, ZH-C44-zuo, ZH-C44-chuzuche

\chapter{Places in Town}
\label{ch:zh-places-in-town}
\hlchaptermodality{44}

\begin{chapteropening}
\textbf{By the end of this chapter:} \emph{I can name places in town and say how I get there.}

Complete the last lesson of chapter 44: I can name places in town and say how I get there.
\end{chapteropening}

\section[\texorpdfstring{\zh{车} (chē)}{车 (chē)}]{\zh{车} (chē) — car; vehicle}

\label{lesson:ZH-C44-che}



\begin{quote}
Before the new one: say the Mandarin for to go out, then the Mandarin for school.
\end{quote}

\subsection*{You'll want to know: \zh{车}}
\textbf{\zh{车}} — \emph{chē} — \textquotedblleft{}car; vehicle\textquotedblright{}.

\begin{quote}
\textbf{\zh{你有车吗}} \emph{nǐ yǒu chē ma} — Have you got a car?
\end{quote}

The new character, in pieces: \zh{车} — four strokes: a bar, a stroke that falls and turns, a long bar, and a vertical through them. The next lesson puts it in your hand.

\begin{cousinweb}
\textbf{\zh{车}} (\emph{chē}) is inherited from Middle Chinese \emph{tsyhae} (Guangyun, in Baxter's transcription), which the dictionary explains as a cart, first made, it says, by the Yellow Emperor.
\end{cousinweb}

\begin{grammarlens}[title={What you've built}]
One more word for everyday answers.
\end{grammarlens}

\subsection*{Your turn}
\begin{itemize}
\raggedright
  \item \emph{Say it:} \emph{chē}
  \item \emph{Say it:} \emph{chē}, once more
  \item \emph{Say it:} \emph{nǐ yǒu chē ma} — Have you got a car?
  \item \emph{Recall:} say the Mandarin for to go out, then the Mandarin for school, then say \emph{chē} again
\end{itemize}

\begin{tcolorbox}[breakable,colback=teal!4,colframe=teal!35!black,title={Before you move on}]
What does \zh{车} mean? (\textquotedblleft{}Car; vehicle\textquotedblright{}.) Say it once more.
\end{tcolorbox}
\section[\texorpdfstring{\zh{车} (chē)}{车 (chē)}]{\zh{车} — four strokes, the vertical last}

\label{lesson:ZH-W44-che}



\begin{quote}
The word first, then the character inside it.

\begin{itemize}
\raggedright
  \item \emph{Recall:} say \emph{chē} — \textbf{R1}, one lesson back, before it decays
\end{itemize}

\begin{quote}
\zh{十} \zh{去}
\end{quote}

\emph{Read it:} each one aloud, with nothing said first — \emph{shí}, \emph{qù}
\end{quote}

\subsection*{Script — four strokes, the vertical last}
\hlblockfigure{\detokenize{figures/ZH-W44-che-filmstrip.pdf}}{How \zh{车} is written, stroke by stroke}

\begin{quote}
\zh{车}
\end{quote}

\textbf{One piece, four strokes.}

\textbf{A short horizontal} at the top.

\textbf{A stroke that falls to the left and turns right}, along under the first bar.

\textbf{A long horizontal} across the bottom.

\textbf{Last, the vertical}, from the top bar down through everything to the foot.

\subsection*{Your turn}
\begin{itemize}
\raggedright
  \item \emph{Write it:} \zh{车}, all four, with the model in front of you
  \item \emph{Cover:} write it again from memory, and count the pieces, not the strokes
\end{itemize}

\begin{tcolorbox}[breakable,colback=teal!4,colframe=teal!35!black,title={Before you move on}]
How many strokes in \zh{车}? (\textbf{Four}.) Which stroke is last? (\textbf{The vertical}.) Which is the longest bar? (\textbf{The third stroke, along the bottom}.)

Source: \href{https://raw.githubusercontent.com/chanind/hanzi-writer-data/68d10a4b21150cae5e1ebbd223eed289cf32d90c/data/\%E8\%BD\%A6.json}{Hanzi Writer Data \zh{车}, pinned stroke paths and medians}.
\end{tcolorbox}
\section[fàndiàn]{fàndiàn — restaurant}

\label{lesson:ZH-C44-fandian}



\begin{quote}
Before the new one: say the Mandarin for school, then the Mandarin for car.
\end{quote}

\subsection*{You'll want to know: fàndiàn}
\emph{fàndiàn} — \textquotedblleft{}restaurant\textquotedblright{}.

\begin{quote}
\emph{wǒmen qù fàndiàn chī fàn} — We're going out to a restaurant.
\end{quote}

Its characters come later; for now it is a word for the ear.

\begin{cousinweb}
\emph{fàndiàn} is built from \zh{饭} 'meal' and diàn 'shop', both inherited from Middle Chinese, \emph{bjonH} and \emph{temH} (Guangyun, in Baxter's transcription).
\end{cousinweb}

\begin{grammarlens}[title={What you've built}]
One more word for everyday answers.
\end{grammarlens}

\subsection*{Your turn}
\begin{itemize}
\raggedright
  \item \emph{Say it:} \emph{fàndiàn}
  \item \emph{Say it:} \emph{fàndiàn}, once more
  \item \emph{Say it:} \emph{wǒmen qù fàndiàn chī fàn} — We're going out to a restaurant.
  \item \emph{Recall:} say the Mandarin for school, then the Mandarin for car, then say \emph{fàndiàn} again
\end{itemize}

\begin{tcolorbox}[breakable,colback=teal!4,colframe=teal!35!black,title={Before you move on}]
What does fàndiàn mean? (\textquotedblleft{}Restaurant\textquotedblright{}.) Say it once more.
\end{tcolorbox}
\section[yīyuàn]{yīyuàn — hospital}

\label{lesson:ZH-C44-yiyuan}



\begin{quote}
Before the new one: say the Mandarin for car, then the Mandarin for restaurant.
\end{quote}

\subsection*{You'll want to know: yīyuàn}
\emph{yīyuàn} — \textquotedblleft{}hospital\textquotedblright{}.

\begin{quote}
\emph{wǒ māma qù yīyuàn le} — My mum has gone to the hospital.
\end{quote}

Its characters come later; for now it is a word for the ear.

\begin{cousinweb}
\emph{yīyuàn} is built from yī 'healing' and yuàn 'courtyard', both inherited from Middle Chinese, \emph{'i} and \emph{hjwenH} (Guangyun, in Baxter's transcription).
\end{cousinweb}

\begin{grammarlens}[title={What you've built}]
One more word for everyday answers.
\end{grammarlens}

\subsection*{Your turn}
\begin{itemize}
\raggedright
  \item \emph{Say it:} \emph{yīyuàn}
  \item \emph{Say it:} \emph{yīyuàn}, once more
  \item \emph{Say it:} \emph{wǒ māma qù yīyuàn le} — My mum has gone to the hospital.
  \item \emph{Recall:} say the Mandarin for car, then the Mandarin for restaurant, then say \emph{yīyuàn} again
\end{itemize}

\begin{tcolorbox}[breakable,colback=teal!4,colframe=teal!35!black,title={Before you move on}]
What does yīyuàn mean? (\textquotedblleft{}Hospital\textquotedblright{}.) Say it once more.
\end{tcolorbox}
\section[zuò]{zuò — to sit; to travel by}

\label{lesson:ZH-C44-zuo}



\begin{quote}
Before the new one: say the Mandarin for restaurant, then the Mandarin for hospital.
\end{quote}

\subsection*{You'll want to know: zuò}
\emph{zuò} — \textquotedblleft{}to sit; to travel by\textquotedblright{}.

\begin{quote}
\emph{wǒmen zuò chē qù} — We're going by car.
\end{quote}

Its characters come later; for now it is a word for the ear.

\begin{cousinweb}
\emph{zuò} is inherited from Middle Chinese \emph{dzwaX} (Guangyun, in Baxter's transcription), which the dictionary explains as bending the joints. Its use for travelling by car or bus is not given there.
\end{cousinweb}

\begin{grammarlens}[title={What you've built}]
One more word for everyday answers.
\end{grammarlens}

\subsection*{Your turn}
\begin{itemize}
\raggedright
  \item \emph{Say it:} \emph{zuò}
  \item \emph{Say it:} \emph{zuò}, once more
  \item \emph{Say it:} \emph{wǒmen zuò chē qù} — We're going by car.
  \item \emph{Recall:} say the Mandarin for restaurant, then the Mandarin for hospital, then say \emph{zuò} again
\end{itemize}

\begin{tcolorbox}[breakable,colback=teal!4,colframe=teal!35!black,title={Before you move on}]
What does zuò mean? (\textquotedblleft{}To sit; to travel by\textquotedblright{}.) Say it once more.
\end{tcolorbox}
\section[chūzūchē]{chūzūchē — taxi}

\label{lesson:ZH-C44-chuzuche}



\begin{quote}
Before the new one: say the Mandarin for hospital, then the Mandarin for to sit.
\end{quote}

\subsection*{You'll want to know: chūzūchē}
\emph{chūzūchē} — \textquotedblleft{}taxi\textquotedblright{}.

\begin{quote}
\emph{chūzūchē hěn guì} — Taxis are expensive.
\end{quote}

Its characters come later; for now it is a word for the ear.

\begin{cousinweb}
\emph{chūzūchē} is built from chūzū 'to rent out' and \zh{车} 'car'; its parts are inherited from Middle Chinese \emph{tsyhwit}, \emph{tsu} 'tax' and \emph{tsyhae} (Guangyun, in Baxter's transcription).
\end{cousinweb}

\begin{grammarlens}[title={What you've built}]
One more word for everyday answers.
\end{grammarlens}

\subsection*{Your turn}
\begin{itemize}
\raggedright
  \item \emph{Say it:} \emph{chūzūchē}
  \item \emph{Say it:} \emph{chūzūchē}, once more
  \item \emph{Say it:} \emph{chūzūchē hěn guì} — Taxis are expensive.
  \item \emph{Recall:} say the Mandarin for hospital, then the Mandarin for to sit, then say \emph{chūzūchē} again
\end{itemize}

\begin{tcolorbox}[breakable,colback=teal!4,colframe=teal!35!black,title={Before you move on}]
What does chūzūchē mean? (\textquotedblleft{}Taxi\textquotedblright{}.) Say it once more.
\end{tcolorbox}
